% ============================================================================
%  ESTADÍSTICA DESCRIPTIVA — Sesión 5
%  Semana 5: 7 al 12 de septiembre de 2026
%  1.8 Aplicación del proceso estadístico en casos prácticos
%  Repaso integrador de la Unidad 1  ·  Preparación del Examen 1
%  Duración prevista: 3 horas
% ============================================================================
\documentclass[aspectratio=169,10pt]{beamer}

\usepackage[utf8]{inputenc}
\IfFileExists{spanish.ldf}{\usepackage[spanish,es-noshorthands]{babel}}{}
\usepackage{amsmath,amssymb}
\usepackage{booktabs}
\usepackage{tikz}
\usepackage{listings}
\usepackage{pgfplots}
\pgfplotsset{compat=1.16}

\usetheme{Uniminuto}

\uniasignatura{Estadística descriptiva}
\uniprofesor{Andrés Felipe Puerta Velez}
\unicorreo{andres.puerta-v@uniminuto.edu.co}
\uninrc{CBASBL031 - CBASBL141}
\unisemana{Sesión 5 — Semana 5 · 7 al 12 de septiembre de 2026}

\newcommand{\listatiny}{\setbeamerfont{itemize/enumerate body}{size=\tiny}\tiny}

\pgfplotsset{
  uniplot/.style={
    width=8.0cm, height=3.7cm,
    axis lines=left,
    tick label style={font=\tiny,color=uniGris},
    label style={font=\tiny,color=uniGris},
    title style={font=\scriptsize\bfseries,color=uniAzul},
    every axis plot/.append style={fill=uniAzul!70,draw=uniAzul},
    ymajorgrids, grid style={uniGris!25,very thin},
  }
}

\newcommand{\fx}[1]{{\ttfamily\color{uniAzul}#1}}

% ------------------------------------------------------------- CÓDIGO R ----
\lstdefinestyle{Rcode}{
  language=R,
  basicstyle=\ttfamily\scriptsize,
  keywordstyle=\color{uniAzul}\bfseries,
  commentstyle=\color{uniGris}\itshape,
  stringstyle=\color{uniTitulo},
  showstringspaces=false, numbers=none, breaklines=true,
  columns=fullflexible,
  backgroundcolor=\color{uniCremaClaro},
  frame=leftline, framerule=1.2pt, rulecolor=\color{uniNaranja},
  xleftmargin=0.35cm, framexleftmargin=0.35cm,
  framexrightmargin=0.15cm, framextopmargin=2pt, framexbottommargin=2pt,
  aboveskip=0.15cm, belowskip=0.15cm
}
\lstdefinestyle{Rmini}{
  language=R,
  basicstyle=\ttfamily\tiny,
  keywordstyle=\color{uniAzul}\bfseries,
  commentstyle=\color{uniGris}\itshape,
  stringstyle=\color{uniTitulo},
  showstringspaces=false, numbers=none, breaklines=true,
  columns=fullflexible,
  backgroundcolor=\color{uniCremaClaro},
  frame=leftline, framerule=1.2pt, rulecolor=\color{uniNaranja},
  xleftmargin=0.35cm, framexleftmargin=0.35cm,
  framexrightmargin=0.15cm, framextopmargin=2pt, framexbottommargin=2pt,
  aboveskip=0.12cm, belowskip=0.12cm
}
\lstdefinestyle{Rout}{
  basicstyle=\ttfamily\tiny\color{black},
  showstringspaces=false, numbers=none, breaklines=true,
  columns=fullflexible,
  backgroundcolor=\color{uniGris!10},
  frame=leftline, framerule=1.2pt, rulecolor=\color{uniGris},
  xleftmargin=0.35cm, framexleftmargin=0.35cm,
  framexrightmargin=0.15cm, framextopmargin=2pt, framexbottommargin=2pt,
  aboveskip=0.12cm, belowskip=0.12cm
}

\begin{document}

\portada

\begin{frame}{Orden del día}
  \begin{enumerate}
    \item Llamado a lista y revisión de los ejercicios propuestos
    \item El proceso estadístico completo, paso a paso
    \item \azul{Caso integrador: Cafetería El Portal} (taller largo)
    \item \azul{Cómo se redacta un informe descriptivo}
    \item Repaso integrador de la Unidad 1
    \item Simulacro del Examen 1 y recomendaciones
  \end{enumerate}

  \vspace{0.25cm}
  {\footnotesize\color{uniGris} Sesión de 3 horas. La próxima semana
  \destacar{no hay clase nueva}: se presenta el Examen 1 (35\%) en Aulas
  Virtuales, presencial en el aula.}
\end{frame}

\begin{frame}{El proceso estadístico completo}
  \scriptsize
  \begin{center}
  \begin{tabular}{@{}rp{4.3cm}p{6.6cm}@{}}
    \toprule
    & \textbf{Paso} & \textbf{Pregunta que se responde} \\
    \midrule
    1 & Definir el problema & ¿Qué decisión quiero tomar con estos datos? \\
    2 & Definir población y muestra & ¿A quién describo? ¿Con cuántos datos? \\
    3 & Identificar las variables & ¿De qué tipo y en qué escala está cada una? \\
    4 & Organizar: tablas & ¿Agrupo o no? ¿Cuántas clases? \\
    5 & Representar: gráficos & ¿Qué gráfico corresponde a esta variable? \\
    6 & Describir: medidas & (Unidad 2) ¿Dónde está el centro y cuánto varía? \\
    7 & Interpretar y decidir & ¿Qué le digo al gerente y qué le recomiendo? \\
    \bottomrule
  \end{tabular}
  \end{center}

  \vspace{0.1cm}
  \begin{alertblock}{Dónde estamos}
    \scriptsize
    Hoy recorremos los pasos \destacar{1 a 5 y el 7}. El paso 6 es la unidad 2.
    Lo importante es que el proceso \azul{siempre es el mismo}, cambie el
    negocio o el dato.
  \end{alertblock}
\end{frame}

\seccion{Caso integrador: Cafetería El Portal}

\begin{frame}{Paso 1 y 2 — el problema y los datos}
  \begin{exampleblock}{La situación}
    \footnotesize
    La \azul{Cafetería El Portal}, en Bello, quiere decidir si abre una segunda
    sede y en qué zona. También quiere saber si le conviene lanzar un combo de
    precio medio. Para eso encuesta a \destacar{45 clientes} durante una
    semana.
  \end{exampleblock}

  \vspace{0.2cm}
  \begin{columns}[T,onlytextwidth]
    \begin{column}{0.48\textwidth}
      \begin{block}{Población y muestra}
        \scriptsize
        \textbf{Población:} todos los clientes de la cafetería.

        \textbf{Muestra:} los 45 encuestados. \quad $n = 45$

        \textbf{Unidad de observación:} un cliente.
      \end{block}
    \end{column}
    \begin{column}{0.48\textwidth}
      \begin{alertblock}{Las dos preguntas del gerente}
        \scriptsize
        (1) ¿De qué zona vienen mis clientes? \\
        (2) ¿Cuánto consume un cliente por visita, y cómo se reparte ese
        consumo?
      \end{alertblock}
    \end{column}
  \end{columns}
\end{frame}

\begin{frame}{Paso 3 — identificar las variables}
  \footnotesize
  \begin{center}
  \begin{tabular}{@{}p{3.1cm}p{3.5cm}p{2.6cm}p{2.6cm}@{}}
    \toprule
    \textbf{Variable} & \textbf{Registro} & \textbf{Tipo} & \textbf{Escala} \\
    \midrule
    Zona de residencia & Centro / Norte / Sur / Otras & Cualitativa nominal
      & Nominal \\
    \addlinespace
    Consumo por visita & Miles de pesos & Cuantitativa continua & Razón \\
    \bottomrule
  \end{tabular}
  \end{center}

  \vspace{0.2cm}
  \begin{alertblock}{Por qué este paso va antes que todo}
    \scriptsize
    De esta tabla salen \destacar{todas} las decisiones siguientes: la zona no
    se agrupa en clases ni se acumula, y el consumo sí se agrupa y sí admite
    ojiva. Equivocarse aquí arrastra el error hasta la conclusión final.
  \end{alertblock}
\end{frame}

\begin{frame}{Los datos crudos}
  \footnotesize
  \textbf{Zona de residencia} (45 respuestas):

  \vspace{0.1cm}
  \begin{beamercolorbox}[wd=\textwidth,sep=0.1cm]{block body}
    \centering\tiny\ttfamily
    C N S C O C N S C N C S O N C S C N O C S N C O C N S C N S C O C S N C N S C O C S N C S
  \end{beamercolorbox}
  {\tiny C = Centro \quad N = Norte \quad S = Sur \quad O = Otras}

  \vspace{0.15cm}
  \textbf{Consumo por visita} (miles de pesos), ya ordenado:

  \vspace{0.1cm}
  \begin{beamercolorbox}[wd=\textwidth,sep=0.1cm]{block body}
    \centering\tiny\ttfamily
    6 \ 7 \ 9 \ 10 \ 11 \ 12 \ 13 \ 14 \ 15 \ 16 \ 16 \ 17 \ 17 \ 17 \ 18 \\
    19 \ 19 \ 20 \ 20 \ 21 \ 21 \ 22 \ 22 \ 22 \ 23 \ 23 \ 23 \ 24 \ 25 \ 25 \\
    26 \ 27 \ 27 \ 28 \ 28 \ 29 \ 29 \ 30 \ 31 \ 32 \ 33 \ 34 \ 37 \ 39 \ 42
  \end{beamercolorbox}

  \vspace{0.1cm}
  {\scriptsize\color{uniGris} Primer reflejo profesional: \azul{ordenar} y
  contar cuántos datos hay. Aquí $n = 45$ en las dos variables, porque cada
  encuestado respondió las dos preguntas.}
\end{frame}

\begin{frame}{Paso 4a — tabla de la variable cualitativa}
  \scriptsize
  \begin{columns}[T,onlytextwidth]
    \begin{column}{0.52\textwidth}
      \centering
      \begin{tabular}{@{}lrrrr@{}}
        \toprule
        \textbf{Zona} & \textbf{$f_i$} & \textbf{$h_i$} & \textbf{$\%_i$} &
        \textbf{Áng.} \\
        \midrule
        Centro & 17 & 0{,}378 & 37{,}8 & $136^\circ$ \\
        Norte  & 11 & 0{,}244 & 24{,}4 & $88^\circ$ \\
        Sur    & 11 & 0{,}244 & 24{,}4 & $88^\circ$ \\
        Otras  &  6 & 0{,}133 & 13{,}3 & $48^\circ$ \\
        \midrule
        \textbf{Total} & \textbf{45} & \textbf{1{,}000} & \textbf{100{,}0} &
        $\mathbf{360^\circ}$ \\
        \bottomrule
      \end{tabular}
    \end{column}
    \begin{column}{0.44\textwidth}
      \begin{exampleblock}{El cálculo}
        \scriptsize
        \[ h_1 = \tfrac{17}{45} = 0{,}378 \]
        \[ \text{áng.}_1 = 0{,}378 \times 360 = 136^\circ \]
        Control: $17+11+11+6 = 45$ \checkmark
      \end{exampleblock}
    \end{column}
  \end{columns}

  \vspace{0.1cm}
  \begin{alertblock}{Sin columnas acumuladas}
    \scriptsize
    La zona es \destacar{nominal}: $F_i$ y $H_i$ no se incluyen. ``El 60\% vive
    en Centro o Norte'' sería un número que depende del orden en que escribí
    las filas.
  \end{alertblock}
\end{frame}

\begin{frame}{Paso 4b — tabla de la variable cuantitativa}
  \scriptsize
  \begin{columns}[T,onlytextwidth]
    \begin{column}{0.40\textwidth}
      \begin{block}{Los cálculos previos}
        \scriptsize
        \[ R = 42 - 6 = 36 \]
        \[ k = 1+3{,}322\log_{10}(45) \]
        \[ k = 1+5{,}49 = 6{,}49 \]
        \[ \Rightarrow k = 6 \]
        \[ A = \tfrac{36}{6} = 6 \]
      \end{block}
      \begin{alertblock}{¿6 o 7 clases?}
        \tiny
        Sturges da 6{,}49: ambos son defendibles. Se toma \destacar{6} porque
        deja una amplitud exacta de 6 y clases redondas. Sturges orienta, no
        obliga.
      \end{alertblock}
    \end{column}
    \begin{column}{0.57\textwidth}
      \centering
      \begin{tabular}{@{}lrrrrrr@{}}
        \toprule
        \textbf{Clase} & \textbf{$x_i$} & \textbf{$f_i$} & \textbf{$F_i$} &
        \textbf{$h_i$} & \textbf{$H_i$} & \textbf{$\%_i$} \\
        \midrule
        $[6,\ 12)$  &  9 &  5 &  5 & 0{,}111 & 0{,}111 & 11{,}1 \\
        $[12,\ 18)$ & 15 &  9 & 14 & 0{,}200 & 0{,}311 & 20{,}0 \\
        $[18,\ 24)$ & 21 & 13 & 27 & 0{,}289 & 0{,}600 & 28{,}9 \\
        $[24,\ 30)$ & 27 & 10 & 37 & 0{,}222 & 0{,}822 & 22{,}2 \\
        $[30,\ 36)$ & 33 &  5 & 42 & 0{,}111 & 0{,}933 & 11{,}1 \\
        $[36,\ 42]$ & 39 &  3 & 45 & 0{,}067 & 1{,}000 &  6{,}7 \\
        \midrule
        \textbf{Tot.} & & \textbf{45} & & \textbf{1{,}000} & &
        \textbf{100{,}0} \\
        \bottomrule
      \end{tabular}
    \end{column}
  \end{columns}
\end{frame}

\begin{frame}{Paso 5 — los gráficos}
  \centering
  \begin{columns}[T,onlytextwidth]
    \begin{column}{0.49\textwidth}
      \begin{tikzpicture}
        \begin{axis}[uniplot, width=6.3cm, height=3.5cm,
          title={Zona de residencia},
          ybar, bar width=0.5cm, ylabel={$f_i$},
          symbolic x coords={Centro,Norte,Sur,Otras},
          xtick=data, ymin=0, ymax=19, enlarge x limits=0.22,
          nodes near coords,
          nodes near coords style={font=\tiny,color=uniAzul}]
          \addplot coordinates {(Centro,17) (Norte,11) (Sur,11) (Otras,6)};
        \end{axis}
      \end{tikzpicture}
    \end{column}
    \begin{column}{0.49\textwidth}
      \begin{tikzpicture}
        \begin{axis}[uniplot, width=6.3cm, height=3.5cm,
          title={Consumo por visita (miles \$)},
          ybar interval, bar width=1, ylabel={$f_i$},
          xtick={6,12,18,24,30,36,42}, ymin=0, ymax=14.5]
          \addplot coordinates {(6,5) (12,9) (18,13) (24,10) (30,5) (36,3)
                                (42,0)};
        \end{axis}
      \end{tikzpicture}
    \end{column}
  \end{columns}

  \vspace{0.05cm}
  \begin{alertblock}{Lo que muestran}
    \scriptsize
    Barras \destacar{separadas} para la zona (nominal, ordenada de mayor a
    menor). Barras \destacar{pegadas} para el consumo (continua). El
    histograma tiene un pico en $[18,24)$ y una \azul{cola a la derecha}:
    distribución asimétrica positiva.
  \end{alertblock}
\end{frame}

\begin{frame}{Paso 5 — la ojiva y lo que permite leer}
  \centering
  \begin{tikzpicture}
    \begin{axis}[uniplot, width=7.6cm, height=3.5cm,
      title={Ojiva del consumo por visita},
      xlabel={Consumo (miles \$)}, ylabel={$F_i$},
      xtick={6,12,18,24,30,36,42}, ymin=0, ymax=49, xmin=5, xmax=43,
      every axis plot/.append style={fill=none}]
      \addplot[color=uniAzul,thick,mark=*,mark size=1.3pt,
               mark options={fill=uniAzul}]
        coordinates {(6,0) (12,5) (18,14) (24,27) (30,37) (36,42) (42,45)};
      \draw[uniTitulo,dashed] (axis cs:5,22.5) -- (axis cs:21.9,22.5)
                              -- (axis cs:21.9,0);
    \end{axis}
  \end{tikzpicture}

  \vspace{-0.05cm}
  \begin{alertblock}{Tres lecturas para el gerente}
    \scriptsize
    (1) $H_3 = 0{,}600$: el \destacar{60\%} consume menos de 24 mil. \quad
    (2) $1-H_4 = 0{,}178$: solo el \destacar{17,8\%} consume 30 mil o más.
    \quad (3) Buscando $F = 22{,}5$ se obtiene un consumo mediano de
    \destacar{unos 22 mil}.
  \end{alertblock}
\end{frame}

\begin{frame}{Paso 7 — de los números a la decisión}
  \begin{block}{Respuesta a la pregunta 1: ¿dónde abrir la segunda sede?}
    \scriptsize
    El \azul{Centro} concentra el 37,8\% de los clientes. Norte y Sur quedan
    \destacar{empatadas} con el 24,4\% cada una (juntas suman el 48,9\%).
    Si la sede actual está en el Centro, Norte y Sur son las candidatas más
    razonables para la segunda sede; el empate obliga a decidir con un
    criterio adicional (costo del local, cercanía, competencia en la zona).
  \end{block}

  \vspace{0.1cm}
  \begin{block}{Respuesta a la pregunta 2: ¿el combo de precio medio?}
    \scriptsize
    El 60\% de los clientes consume menos de 24 mil y la clase más frecuente
    es \azul{$[18, 24)$}. Un combo alrededor de \destacar{20 a 22 mil} apunta
    justo al grupo más numeroso. Los clientes de 30 mil o más son solo el
    17,8\%: no justifican una línea premium por ahora.
  \end{block}

  \vspace{0.1cm}
  \begin{alertblock}{El límite de lo que podemos afirmar}
    \scriptsize
    Todo esto describe a los \destacar{45 encuestados}. Extenderlo a todos los
    clientes es hablar de quienes no se midieron. Además, la muestra se tomó
    \azul{una sola semana}:
    si esa semana fue atípica, las conclusiones cambian.
  \end{alertblock}
\end{frame}

\seccion{Cómo se redacta el informe}

\begin{frame}{Plantilla de informe descriptivo}
  \scriptsize
  \begin{center}
  \begin{tabular}{@{}p{2.9cm}p{9.0cm}@{}}
    \toprule
    \textbf{Sección} & \textbf{Qué va ahí} \\
    \midrule
    Objetivo & Una frase: qué se quería saber y para qué decisión \\
    Datos y método & Población, muestra, $n$, cómo se recogieron, fecha \\
    Variables & Tabla con tipo y escala de cada una \\
    Resultados & Tablas y gráficos, cada uno con título y numeración \\
    Interpretación & Lo que dicen los números, en lenguaje del negocio \\
    Limitaciones & Tamaño de muestra, periodo, sesgos posibles \\
    Recomendación & Qué se propone hacer, y con qué evidencia \\
    \bottomrule
  \end{tabular}
  \end{center}

  \vspace{0.1cm}
  \begin{alertblock}{La regla de oro de la redacción}
    \scriptsize
    Cada cifra que aparezca debe venir acompañada de \destacar{qué mide, sobre
    cuántos casos y en qué unidad}. ``El promedio es 21'' no dice nada;
    ``el consumo promedio de los 45 encuestados fue de 21 mil pesos por
    visita'' sí.
  \end{alertblock}
\end{frame}

\begin{frame}{Ejercicio en clase — escriba el informe}
  \begin{alertblock}{En grupos de 3 (25 minutos)}
    \footnotesize
    Con las tablas y gráficos de la Cafetería El Portal, redacte un informe de
    \destacar{máximo una página} siguiendo la plantilla.

    \vspace{0.2cm}
    \begin{enumerate}\setlength{\itemsep}{-0.05cm}
      \item Escriba el objetivo en una sola frase.
      \item Describa datos y método en tres líneas.
      \item Redacte tres hallazgos, cada uno con su cifra y su unidad.
      \item Enuncie dos limitaciones reales del estudio.
      \item Cierre con una recomendación concreta y accionable.
    \end{enumerate}

    \vspace{0.15cm}
    Al final, dos grupos leen su informe y el curso señala qué cifras están
    mal enunciadas.
  \end{alertblock}
\end{frame}

\seccion{Repaso integrador de la Unidad 1}

\begin{frame}{Mapa de la unidad}
  \centering
  \begin{tikzpicture}[scale=0.92,every node/.style={font=\tiny,align=center},
    nodo/.style={draw=uniAzul,fill=uniAzul!10,rounded corners,inner sep=3pt,
                 minimum width=2.3cm,minimum height=0.6cm},
    flecha/.style={->,>=latex,uniGris}]
    \node[nodo] (a) at (0,0) {Población y muestra\\ parámetro / estadístico};
    \node[nodo] (b) at (3.6,0) {Tipo de variable\\ y escala};
    \node[nodo] (c) at (7.2,0) {Tabla de\\ frecuencias};
    \node[nodo] (d) at (10.8,0) {Gráfico};
    \node[nodo,fill=uniCrema,draw=uniTitulo] (e) at (7.2,-1.6)
      {Interpretación y decisión};
    \node[nodo,fill=uniCremaClaro,draw=uniNaranja] (f) at (3.6,-1.6)
      {¿Agrupada o\\ no agrupada?};
    \draw[flecha] (a)--(b); \draw[flecha] (b)--(c); \draw[flecha] (c)--(d);
    \draw[flecha] (b)--(f); \draw[flecha] (f)--(c);
    \draw[flecha] (c)--(e); \draw[flecha] (d)--(e);
  \end{tikzpicture}

  \vspace{0.2cm}
  \begin{alertblock}{Si tuviera que recordar una sola idea}
    \scriptsize
    \destacar{La escala de la variable manda.} Determina si se agrupa, si se
    acumula, qué gráfico corresponde y —desde la unidad 2— qué medida se
    puede calcular.
  \end{alertblock}
\end{frame}

\begin{frame}{Formulario de la Unidad 1}
  \scriptsize
  \begin{columns}[T,onlytextwidth]
    \begin{column}{0.48\textwidth}
      \begin{block}{Frecuencias}
        \scriptsize
        \[ h_i = \frac{f_i}{n} \qquad \%_i = h_i \times 100 \]
        \[ F_i = \sum_{j\le i} f_j \qquad H_i = \sum_{j\le i} h_j \]
        \[ \sum f_i = n \qquad \sum h_i = 1 \]
      \end{block}
    \end{column}
    \begin{column}{0.48\textwidth}
      \begin{block}{Agrupamiento}
        \scriptsize
        \[ R = x_{\max} - x_{\min} \]
        \[ k = 1 + 3{,}322\,\log_{10}(n) \]
        \[ A = \frac{R}{k} \ \text{(hacia arriba)} \]
        \[ x_i = \frac{L_{\inf}+L_{\sup}}{2} \]
      \end{block}
    \end{column}
  \end{columns}

  \vspace{0.1cm}
  \begin{exampleblock}{Sector circular y lectura de acumuladas}
    \scriptsize
    \[ \text{ángulo}_i = h_i \times 360^\circ \qquad
       P(X < L_{\sup,i}) = H_i \qquad P(X \ge L_{\sup,i}) = 1 - H_i \]
  \end{exampleblock}
\end{frame}

\begin{frame}{Simulacro del Examen 1 — parte A}
  \begin{alertblock}{Individual (15 minutos, sin apuntes)}
    \scriptsize
    \begin{enumerate}\setlength{\itemsep}{-0.05cm}
      \item Una empresa mide el tiempo de entrega de sus 800 pedidos y toma 60.
            Identifique población, muestra y $n$.
      \item ``El tiempo promedio de los 60 pedidos es 3,2 días.''
            ¿Parámetro o estadístico? ¿Símbolo?
      \item Clasifique: código de producto; nivel de riesgo (bajo/medio/alto);
            unidades vendidas; peso del paquete.
      \item ¿Cuáles de las cuatro admiten media? ¿Cuáles admiten moda?
      \item Con $n = 60$, $x_{\min} = 1$ y $x_{\max} = 13$: calcule $R$, $k$ y $A$.
      \item Escriba los intervalos de clase resultantes.
    \end{enumerate}
  \end{alertblock}
\end{frame}

\begin{frame}{Simulacro — solución de la parte A}
  \scriptsize
  \begin{columns}[T,onlytextwidth]
    \begin{column}{0.49\textwidth}
      \begin{block}{Preguntas 1 a 4}
        \scriptsize
        \textbf{1.} Población: los 800 pedidos. Muestra: los 60 medidos.
        $n = 60$.

        \textbf{2.} Estadístico: $\bar{x} = 3{,}2$.

        \textbf{3.} Código: cualitativa nominal. Riesgo: cualitativa ordinal.
        Unidades: cuantitativa discreta. Peso: cuantitativa continua.

        \textbf{4.} Media: unidades y peso. Moda: las cuatro.
      \end{block}
    \end{column}
    \begin{column}{0.49\textwidth}
      \begin{exampleblock}{Preguntas 5 y 6}
        \scriptsize
        \[ R = 13-1 = 12 \]
        \[ k = 1+3{,}322\log_{10}(60) = 1+5{,}91 = 6{,}91 \Rightarrow 7 \]
        \[ A = \tfrac{12}{7} = 1{,}71 \Rightarrow 2 \]
        Intervalos: $[1,3)$, $[3,5)$, $[5,7)$, $[7,9)$, $[9,11)$, $[11,13)$,
        $[13,15]$.
      \end{exampleblock}
    \end{column}
  \end{columns}

  \vspace{0.05cm}
  {\tiny\color{uniGris} Nota: al redondear $A$ hacia arriba, $k\cdot A = 14 > R$;
  la última clase puede quedar vacía o puede reducirse a 6 clases con $A=2$
  ajustando el inicio. Ambas soluciones se aceptan si están justificadas.}
\end{frame}

\begin{frame}{Simulacro del Examen 1 — parte B}
  \begin{alertblock}{Individual (15 minutos)}
    \scriptsize
    Se tiene la siguiente tabla de la variable ``ventas diarias'' (millones):

    \vspace{0.15cm}
    \centering
    \begin{tabular}{@{}lrrrrr@{}}
      \toprule
      \textbf{Clase} & $[10,20)$ & $[20,30)$ & $[30,40)$ & $[40,50)$ &
      $[50,60]$ \\
      \midrule
      $f_i$ & 6 & 14 & 20 & 8 & 2 \\
      \bottomrule
    \end{tabular}

    \vspace{0.2cm}
    \raggedright
    \begin{enumerate}\setlength{\itemsep}{-0.05cm}
      \item Complete $x_i$, $F_i$, $h_i$ y $H_i$.
      \item ¿Qué porcentaje de días vendió menos de 40 millones?
      \item ¿Qué porcentaje vendió 40 millones o más?
      \item ¿Cuál es la clase modal?
      \item ¿Qué forma tiene la distribución? Justifique sin calcular nada.
      \item ¿Qué gráfico usaría para mostrarle esto a la junta y por qué?
    \end{enumerate}
  \end{alertblock}
\end{frame}

\begin{frame}{Simulacro — solución de la parte B}
  \scriptsize
  \begin{columns}[T,onlytextwidth]
    \begin{column}{0.50\textwidth}
      \centering
      \begin{tabular}{@{}lrrrrr@{}}
        \toprule
        \textbf{Clase} & \textbf{$x_i$} & \textbf{$f_i$} & \textbf{$F_i$} &
        \textbf{$h_i$} & \textbf{$H_i$} \\
        \midrule
        $[10,20)$ & 15 &  6 &  6 & 0{,}12 & 0{,}12 \\
        $[20,30)$ & 25 & 14 & 20 & 0{,}28 & 0{,}40 \\
        $[30,40)$ & 35 & 20 & 40 & 0{,}40 & 0{,}80 \\
        $[40,50)$ & 45 &  8 & 48 & 0{,}16 & 0{,}96 \\
        $[50,60]$ & 55 &  2 & 50 & 0{,}04 & 1{,}00 \\
        \midrule
        \textbf{Tot.} & & \textbf{50} & & \textbf{1{,}00} & \\
        \bottomrule
      \end{tabular}
    \end{column}
    \begin{column}{0.47\textwidth}
      \begin{exampleblock}{Respuestas 2 a 6}
        \scriptsize
        \textbf{2.} $H_3 = 0{,}80 \Rightarrow \destacar{80\%}$

        \textbf{3.} $1-0{,}80 = \destacar{20\%}$

        \textbf{4.} Clase modal: $[30,40)$, con $f=20$

        \textbf{5.} Asimétrica \azul{negativa}: el pico está a la derecha del
        centro y la cola larga queda a la izquierda.

        \textbf{6.} Histograma: la variable es continua y agrupada, y la junta
        necesita ver la forma.
      \end{exampleblock}
    \end{column}
  \end{columns}
\end{frame}

\begin{frame}{Recomendaciones para el Examen 1}
  \begin{block}{Qué se evalúa}
    \scriptsize
    \begin{itemize}\setlength{\itemsep}{0.02cm}
      \item Identificar población, muestra, parámetro y estadístico.
      \item Clasificar variables y escalas, y decir qué se puede calcular.
      \item Construir tablas de frecuencia, agrupadas y no agrupadas.
      \item Leer e interpretar $F_i$ y $H_i$ en lenguaje corriente.
      \item Elegir el gráfico correcto y detectar gráficos engañosos.
    \end{itemize}
  \end{block}

  \vspace{0.1cm}
  \begin{alertblock}{Cinco consejos concretos}
    \scriptsize
    (1) Lea el enunciado dos veces y subraye ``a lo sumo'', ``menos de'',
    ``más de''. \quad (2) Verifique siempre $\sum f_i = n$. \quad
    (3) Escriba las unidades. \quad (4) Redondee al final, no en cada paso.
    \quad (5) Si le piden interpretar, \destacar{escriba una frase completa},
    no un número suelto.
  \end{alertblock}

  \vspace{0.05cm}
  {\tiny\color{uniGris} Examen 1: semana 6 (14--19 de septiembre), 35\% de la
  nota, en Aulas Virtuales (Moodle), presencial en el aula.}
\end{frame}

\begin{frame}[fragile]{El caso completo, en R}
  \scriptsize
  Todo el proceso estadístico de hoy —pasos 3 a 5— en un solo guion:
\begin{lstlisting}[style=Rmini]
# Paso 3: declarar el tipo de cada variable
zona    <- factor(zona)                 # nominal
consumo <- as.numeric(consumo)          # continua

# Paso 4a: tabla de la cualitativa
table(zona)
round(prop.table(table(zona)) * 100, 1)

# Paso 4b: tabla de la cuantitativa
limites <- seq(6, 42, by = 6)
clases  <- cut(consumo, breaks = limites, right = FALSE, include.lowest = TRUE)
data.frame(xi = head(limites, -1) + 3,
           fi = as.vector(table(clases)),
           Fi = as.vector(cumsum(table(clases))),
           hi = round(as.vector(prop.table(table(clases))), 3))

# Paso 5: los graficos
barplot(table(zona)); hist(consumo, breaks = limites, right = FALSE)
\end{lstlisting}

  \vspace{-0.1cm}
  {\tiny\color{uniGris} La ventaja no es la velocidad: es que el guion
  \destacar{queda escrito}. Si mañana llegan 20 encuestas más, se cambia el
  vector y todo se recalcula solo.}
\end{frame}

\begin{frame}{Ejercicios propuestos}
  \vspace{-0.1cm}
  \begin{beamercolorbox}[wd=\textwidth,sep=0.10cm]{block body}
    \centering\scriptsize
    Repaso general de la Unidad 1. Resuélvalos antes del examen.
  \end{beamercolorbox}
  \setbeamerfont{itemize/enumerate body}{size=\tiny}\tiny
  \begin{columns}[T,onlytextwidth]
    \begin{column}{0.47\textwidth}
      \begin{enumerate}\setlength{\itemsep}{0.16cm}
        \item Defina población, muestra y unidad de observación
        \item Diferencia entre parámetro y estadístico, con símbolos
        \item Dé un ejemplo propio de cada tipo de variable
        \item ¿Por qué el estrato no admite media?
        \item ¿Qué distingue la escala intervalar de la de razón?
        \item Con $n=45$, $f_2=9$: calcule $h_2$ y $\%_2$
        \item Si $H_3=0{,}60$: ¿qué significa exactamente?
        \item $x_{\min}=6$, $x_{\max}=42$, $n=45$: calcule $R$, $k$, $A$
        \item Escriba los intervalos del ejercicio anterior
        \item Marca de clase de $[24,30)$
      \end{enumerate}
    \end{column}
    \begin{column}{0.50\textwidth}
      \begin{enumerate}\setlength{\itemsep}{0.16cm}
        \setcounter{enumi}{10}
        \item ¿Cuál es la clase modal del consumo en El Portal?
        \item ¿Qué \% consume menos de 18 mil?
        \item ¿Qué \% consume 24 mil o más?
        \item Ángulo del sector ``Norte'' en el circular
        \item ¿Por qué la zona no lleva $F_i$?
        \item ¿Qué gráfico para la zona? ¿Y para el consumo?
        \item Enuncie los 7 pasos del proceso estadístico
        \item Escriba dos limitaciones del estudio de El Portal
        \item Redacte una recomendación de una línea para el gerente
        \item ¿Por qué no se puede afirmar algo sobre quien no se midió?
      \end{enumerate}
    \end{column}
  \end{columns}
\end{frame}

\begin{frame}{Respuestas}{Autoevaluación de los ejercicios propuestos}
  \setbeamerfont{itemize/enumerate body}{size=\tiny}\tiny
  \begin{columns}[T,onlytextwidth]
    \begin{column}{0.47\textwidth}
      \begin{enumerate}\setlength{\itemsep}{0.16cm}
        \item Conjunto total / subconjunto observado / cada elemento
        \item Parámetro describe la población ($\mu,\sigma,p$); estadístico
              la muestra ($\bar x,s,\hat p$)
        \item Respuesta abierta
        \item Es ordinal: las distancias entre estratos no son iguales
        \item El cero de razón es absoluto; el de intervalo es convencional
        \item $h_2 = 0{,}200$; \ $\%_2 = 20{,}0\%$
        \item El 60\% de los datos está por debajo de ese límite superior
        \item $R=36$; \ $k=6{,}49 \Rightarrow 6$; \ $A=6$
        \item $[6,12)$, $[12,18)$, $[18,24)$, $[24,30)$, $[30,36)$, $[36,42]$
        \item $x_i = 27$
      \end{enumerate}
    \end{column}
    \begin{column}{0.50\textwidth}
      \begin{enumerate}\setlength{\itemsep}{0.16cm}
        \setcounter{enumi}{10}
        \item $[18,24)$, con $f_3 = 13$
        \item $H_2 = 31{,}1\%$
        \item $1-H_3 = 40{,}0\%$
        \item $0{,}244\times360 = 88^\circ$
        \item Es nominal: acumular no tendría significado
        \item Barras (o circular) y histograma, respectivamente
        \item Problema, población/muestra, variables, tablas, gráficos,
              medidas, interpretación
        \item Muestra pequeña; una sola semana de observación
        \item Respuesta abierta, con cifra y unidad
        \item Porque los datos solo respaldan lo que se observó; del resto
              no hay evidencia
      \end{enumerate}
    \end{column}
  \end{columns}
\end{frame}

\begin{frame}{Referencias}
  \footnotesize
  \begin{itemize}
    \item[{[1]}] Monroy Saldívar, S. (2008). \textit{Estadística descriptiva}.
      Instituto Politécnico Nacional.
      \texttt{https://elibro.net/es/lc/uniminuto/titulos/74722}
    \item[{[2]}] Ordóñez Fernández, F. F. \& González Fernández, J. (2021).
      \textit{Estadística descriptiva paso a paso}. Ediciones Pirámide.
      \texttt{https://elibro.net/es/lc/uniminuto/titulos/215449}
    \item[{[3]}] Llinás Solano, H. (2014). \textit{Introducción a la
      estadística con aplicaciones en ciencias sociales}. Universidad del
      Norte. \texttt{https://elibro.net/es/lc/uniminuto/titulos/70062}
  \end{itemize}

  \vspace{0.35cm}
  \begin{center}
    {\Large\color{uniAzul}\textbf{Éxitos en el Examen 1}}
  \end{center}
\end{frame}

\end{document}
